\section{Conclusiones y recomendaciones}

\subsection*{Respuesta al objetivo general}

El TFM diseña y evalúa una arquitectura escalonada con componentes distribuidos y cierres globales auditados para formar coaliciones multi-AMR bajo requisitos de cardinalidad, capacidad, wrench, dinámica, seguridad e información. A0--FULL ejecuta los seis prefijos de comprobación sobre la misma carga Euler--Lagrange y los mismos mundos, mientras SP1--SP8 conservan la función de aislar constructos y comparadores.

El objetivo general se cumple dentro del simulador numérico declarado. La campaña integrada correctiva cerró 400 mundos y 2.400 ejecuciones de método, con $n=100$ fijo por familia, cero errores numéricos y semillas abiertas después del \emph{freeze}. Los resultados demuestran que la validez física no es monótona al añadir una prueba: A3 resolvió complementariedad de momento, pero degradó escasez al eliminar reserva dinámica. Este hallazgo vincula decisión y mecánica de una forma que los estudios separados no podían identificar.
\subsection*{Cumplimiento de los objetivos específicos}

Los objetivos se cumplen con alcances distintos. O1 y O2 quedan cubiertos por la definición de coalición física y la escalera integrada y modular. O3 se cumple solo de forma parcial: Smith--QR incorpora memoria y ocupación, pero A0 usa un umbral global; A1--A4 ejecutan ranking, capacidad, búsqueda de contactos o unión con información global; y FULL consulta todos los ociosos. O4 queda sustentado por los juegos de capacidad y contacto de SP2--SP3 dentro del modelo planar. O5 se cumple para comparadores clásicos, centralizados y basados en modelo, y solo de forma documental para MARL: MAPPO--GNN conserva arquitectura, presupuesto, semillas y hashes, pero su gate de sensibilidad falla y bloquea una afirmación de rendimiento aprendido. O6 queda trazado mediante configuraciones congeladas, semillas, hashes, manifiestos, un registro afirmación--artefacto y 25 estimandos auditables.
\subsection*{Aportes teóricos y metodológicos}

La \emph{coalición física} separa partición combinatoria y ejecución. La aplicación teórica derivada formaliza esa relación en el lazo de carga: si el cambio de coalición modifica un error de realización de wrench acotado, no reinicia el estado y conserva la misma matriz cerrada, una solución común de Lyapunov proporciona ISS práctica para cualquier secuencia medible de sustituciones. El teorema no cubre cambios de masa, resets ni contacto tridimensional. V1--V3 dejaron de comparar una fórmula consigo misma: usan una solución KKT independiente, la identidad de potencia por diferencias finitas, un QP SLSQP y una integración lineal separada.

El aporte metodológico combina mundos emparejados, semántica explícita de aceptación y tamaño confirmatorio fijo. A0--FULL distingue \texttt{accepted\_by\_stage}, falso positivo físico y éxito final. Los veinte contrastes fijan endpoint, estimando, unidad independiente, McNemar, \emph{bootstrap}, semilla de análisis y familia Holm antes de abrir test. La sensibilidad SP4 agrega seis escenarios dentro de cada uno de 18 bloques semilla--flota y evita interpretar 108 instancias correlacionadas como 108 réplicas IID. Todos los fallos permanecen en el denominador.
\subsection*{Lectura de las hipótesis}

Las campañas producen resultados favorables, parciales, negativos y no estimables:
La hip\'otesis central recibe apoyo espec\'ifico, no mon\'otono. En escasez, A1 y A2 mejoraron la tasa de \'exito en $0{,}56$ y $0{,}44$, pero A3 la redujo en $0{,}16$: cerrar el \emph{wrench} est\'atico retir\'o reserva de tracci\'on. En momento, obst\'aculo y \emph{dropout}, A3, A4 y FULL aportaron $0{,}63$, $0{,}22$ y $0{,}67$, respectivamente. El \'ultimo efecto pertenece al paquete comunicaci\'on--reemplazo y no permite atribuci\'on individual.

H1--H3 explican la formaci\'on de la coalici\'on, pero no demuestran un ganador universal. Smith+QR redujo el \emph{regret} frente a \emph{greedy}, aunque qued\'o por detr\'as de la subasta proxy y de uniforme+QR. Los \emph{fitness} marginales mejoraron las versiones \emph{plain}, pero uniforme+\emph{closure} obtuvo el menor \emph{regret}. La guardia de \emph{wrench} elimin\'o los falsos positivos condicionados; como uniforme+guardia tambi\'en alcanz\'o cobertura completa, la mejora CLOSED se atribuye principalmente a la reparaci\'on global, no a superioridad distribuida extremo a extremo.

H4--H6 muestran qué ocurre una vez formada la coalición. En SP4 v3, Replicator+CBF obtuvo $0{,}2685$ de éxito seguro y cero colisiones barridas, pero $0{,}7315$ de timeouts; su diferencia frente a CBF fue $0{,}0926$. La sensibilidad sobre 18 bloques arroja 7 bloques favorables, 1 adverso y 10 empates ($p_{\rm Holm}=0{,}0352$), y las otras cuatro direcciones también sobreviven Holm. SP5 v2 separó RAW, SAFE y EXEC sin reparar poses: CBF local alcanzó $0{,}593$ de éxito seguro y cero colisiones; Hamiltoniano+CBF obtuvo $0{,}426$ y cero colisiones, frente a $0{,}167$ de éxito y $0{,}833$ de colisión en RAW. El proxy VO logró $0{,}657$, pero con $0{,}343$ de colisiones, y la referencia centralizada solo $0{,}167$. En SP6, la guardia \emph{wrench-market} redujo la pérdida de carga en $0{,}0227$ frente a \emph{greedy}. Estos resultados pertenecen al simulador planar y no equivalen a seguridad funcional ni a contacto realista.

H7 y H8 no son confirmatorias. La asociaci\'on entre radio y conectividad ($\rho=0{,}8934$) es descriptiva, H7.3 carece de bloques completos y la campa\~na mesosc\'opica no midi\'o de forma externa \emph{timeouts} ni memoria residente. La conclusi\'on correcta es conservar esas preguntas abiertas, no convertir tendencias en evidencia.

\subsection*{Limitaciones}

La evidencia se limita a simulación planar. El agarre se comprueba con fuerzas estáticas unilaterales, mientras el seguimiento supone un agarre rígido capaz de transmitir tracción acotada de ambos signos. Esa hipótesis permite frenar, pero no reproduce contacto, deformación, fricción ni saturación de rueda tridimensionales. El waypoint del caso de obstáculo es nominal; la seguridad se decide por el despeje observado. FULL conserva un 11\% de colisiones en el régimen adverso porque el \emph{wrench} realizado no siempre reproduce la aceleración HOCBF. Por tanto, no se demuestra seguridad funcional.

La red integrada incluye colas, \emph{timestamps} implícitos, pérdida y retardo, pero usa un modelo sintético y un solo \emph{dropout}. El selector de reemplazo consulta todos los robots ociosos y no restringe candidatos por adyacencia; FULL es un proxy de recuperación y no una implementación local extremo a extremo. SP7 histórico sigue siendo descriptivo y SP8 mesoscópico no acredita \emph{timeouts} ni RSS observados. El cierre A3 explora un vecindario acotado de uno o dos robots; no optimiza globalmente carga, contacto y trayectoria. La matriz ISS común exige carga y ganancias fijas y no cubre resets o dinámica de contacto cambiante. El umbral $\rho\leq0{,}16$, las ganancias HOCBF, el paso y el horizonte fueron congelados, pero no se sometieron a sensibilidad paramétrica. Las frecuencias por familia describen el generador y no estiman prevalencia industrial.
\subsection*{Trabajo futuro prioritario}

La progresión futura se gobierna por gates; ninguna etapa hereda automáticamente la
validez de la anterior.

\begin{table}[H]
\centering
\scriptsize
\begin{tabularx}{\textwidth}{@{}p{0.12\textwidth}p{0.23\textwidth}p{0.27\textwidth}X@{}}
\toprule
Gate & Entrada mínima & Criterio de salida & Peligros y responsabilidad \\
\midrule
G1: motor dinámico & Configuración y mundos congelados; controlador cerrado, no replay. & Contacto con fricción/deslizamiento, saturación de rueda, SAFE no realizable y abstención medidos sin reparar trayectorias; hashes y fallos preservados. & Caída de carga y colisión todavía simuladas; responsable: equipo de modelado y control. \\
G2: HIL & G1 cerrado; interfaces temporizadas, emulación de sensores/actuadores y E-stop verificable. & Latencia, pérdida, ruido, saturación, watchdog, parada y recuperación dentro de límites predeclarados. & Comandos obsoletos, pérdida de comunicación y actuación inesperada; responsables: control, integración y seguridad. \\
G3: piloto AMR & G2 cerrado; evaluación de riesgos, zona segregada, operador formado y protocolo de rescate. & Pruebas escalonadas con carga limitada, registro de incidentes, parada manual/automática y criterios de aborto aprobados. & Atrapamiento, impacto, caída de carga, pérdida de red y recuperación manual; responsables: investigador principal, operador y titular de la instalación. \\
\bottomrule
\end{tabularx}
\caption{Ruta propuesta de simulación a motor dinámico, HIL y piloto físico. Es un plan futuro, no evidencia obtenida.}
\label{tab:gates-validacion-futura}
\end{table}

En paralelo, deben incorporarse el margen de autoridad al cierre sin observar test,
extender ISS a matrices dependientes de coalición con Lyapunov múltiple, y medir RSS y
\emph{deadlines} mediante un \emph{watchdog} externo común.
La conclusión central es mecánica y diagnóstica. Una coalición combinatoriamente válida puede fallar por capacidad, contacto, dinámica, seguridad o información; incluso un cierre que mejora el wrench estático puede empeorar el transporte al retirar reserva de tracción. A0--FULL localiza ese intercambio, mientras la cota ISS delimita cuándo la sustitución permanece estable. El resultado es una arquitectura auditable, no una promesa de superioridad universal.
